\chapter{2013年四川卷（理）}

\section{单选题}

\begin{problem}
设集合 \(A = \{x \mid x + 2 = 0\}\)，集合 \(B = \{x \mid x^2 - 4 = 0\}\)，则 \(A \cap B =\)（\quad）

\choices
    {\(\{-2\}\)}
    {\(\{2\}\)}
    {\(\{-2,2\}\)}
    {\(\varnothing\)}
\end{problem}

\begin{answer}
A
\end{answer}

\begin{solution}
由 \(x+2=0\) 得 \(A=\{-2\}\)，由 \(x^2-4=0\) 得 \(x=\pm2\)，所以 \(B=\{-2,2\}\). 因而 \(A\cap B=\{-2\}\)，故选 A.
\end{solution}

\begin{problem}
如图，在复平面内，点 \(A\) 表示复数 \(z\) 的共轭复数，则图中表示复数 \(z\) 的点是（\quad）

\bitmapfigure[max width=0.45\linewidth]{img/2013/sichuan_science/q2_fig1.png}

\choices
    {\(A\)}
    {\(B\)}
    {\(C\)}
    {\(D\)}
\end{problem}

\begin{answer}
B
\end{answer}

\begin{solution}
设 \(z=u+vi\)，则 \(\overline z=u-vi\)，所以表示 \(z\) 与 \(\overline z\) 的点关于 \(x\) 轴对称. 图中点 \(A\) 表示 \(\overline z\)，其关于 \(x\) 轴的对称点是 \(B\)，故选 B.
\end{solution}

\begin{problem}
一个几何体的三视图如图所示，则该几何体的直观图可以是（\quad）

\bitmapfigure[max width=0.55\linewidth]{img/2013/sichuan_science/q3_fig1.png}

\choices
    {\choicebitmap[width=0.15\paperwidth]{img/2013/sichuan_science/q3_optA.png}}
    {\choicebitmap[width=0.15\paperwidth]{img/2013/sichuan_science/q3_optB.png}}
    {\choicebitmap[width=0.15\paperwidth]{img/2013/sichuan_science/q3_optC.png}}
    {\choicebitmap[width=0.15\paperwidth]{img/2013/sichuan_science/q3_optD.png}}
\end{problem}

\begin{answer}
D
\end{answer}

\begin{solution}
俯视图是一个外圆和一个内虚圆，说明该几何体的上、下相应截面均为圆，且下部圆面被上部遮挡. 选项 A、C 的俯视图含有矩形轮廓；选项 B 的下部为棱柱，俯视图的内轮廓不是圆. 只有选项 D 是上部圆台与下部圆柱的组合，正视图、侧视图分别为梯形上接矩形，俯视图为外圆内虚圆，故选 D.
\end{solution}

\begin{problem}
设 \(x \in \Z\)，集合 \(A\) 是奇数集，集合 \(B\) 是偶数集. 若命题 \(p\)：\(\forall x \in A\)，\(2x \in B\)，则（\quad）

\choices
    {\(\neg p\)：\(\forall x \in A\)，\(2x \notin B\)}
    {\(\neg p\)：\(\forall x \notin A\)，\(2x \notin B\)}
    {\(\neg p\)：\(\exists x \notin A\)，\(2x \in B\)}
    {\(\neg p\)：\(\exists x \in A\)，\(2x \notin B\)}
\end{problem}

\begin{answer}
D
\end{answer}

\begin{solution}
命题 \(p\) 是“对所有 \(x\in A\)，都有 \(2x\in B\)”. 全称命题的否定是存在一个反例，因此
\(\neg p\) 为“存在 \(x\in A\)，使 \(2x\notin B\)”，对应选项 D.
\end{solution}

\begin{problem}
函数 \( f(x) = 2\sin\left(\omega x + \varphi\right) \)（\(\omega > 0\)，\(-\frac{\pi}{2} < \varphi < \frac{\pi}{2}\)）的部分图象如图所示，则 \(\omega\)，\(\varphi\) 的值分别是（\quad）

\bitmapfigure[max width=0.45\linewidth]{img/2013/sichuan_science/q5_fig1.png}

\choices
    {\(2\)，\(-\frac{\pi}{3}\)}
    {\(2\)，\(-\frac{\pi}{6}\)}
    {\(4\)，\(-\frac{\pi}{6}\)}
    {\(4\)，\(\frac{\pi}{3}\)}
\end{problem}

\begin{answer}
A
\end{answer}

\begin{solution}
由图象可知最大值 \(2\) 在 \(x=\frac{5\pi}{12}\) 处取得，且在 \(x=-\frac{\pi}{3}\) 处过 \(x\) 轴并由正变负. 因而在这两个点之间，相位从 \(\pi\) 增加到 \(\frac{5\pi}{2}\)，故
\(\omega\left(\frac{5\pi}{12}+\frac{\pi}{3}\right)=\frac{3\pi}{2}\)，\(\omega=2\).
又 \(\omega\frac{5\pi}{12}+\varphi=\frac{\pi}{2}\)，所以
\[
\varphi=\frac{\pi}{2}-2\cdot\frac{5\pi}{12}=-\frac{\pi}{3}
\]
满足给定范围，故选 A.
\end{solution}

\begin{problem}
抛物线 \( y^{2} = 4x \) 的焦点到双曲线 \( x^{2} - \frac{y^{2}}{3} = 1 \) 的渐近线的距离是（\quad）

\choices
    {\( \frac{1}{2} \)}
    {\(\frac{\sqrt{3}}{2}\)}
    {\(1\)}
    {\(\sqrt{3}\)}
\end{problem}

\begin{answer}
B
\end{answer}

\begin{solution}
\(y^2=4x\) 是标准抛物线 \(y^2=4ax\)，故 \(a=1\)，焦点为 \(F(1,0)\). 双曲线
\(x^2-\frac{y^2}{3}=1\) 的渐近线为 \(y=\pm\sqrt3x\). 点 \(F\) 到直线 \(\sqrt3x-y=0\) 的距离为
\[
\frac{|\sqrt3|}{\sqrt{(\sqrt3)^2+(-1)^2}}=\frac{\sqrt3}{2}
\]
两条渐近线关于 \(x\) 轴对称，距离相同，故选 B.
\end{solution}

\begin{problem}
函数 \( y = \frac{x^3}{3^{x} - 1} \) 的图象大致是（\quad）

\choices
    {\choicebitmap[width=0.15\paperwidth]{img/2013/sichuan_science/q7_optA.png}}
    {\choicebitmap[width=0.15\paperwidth]{img/2013/sichuan_science/q7_optB.png}}
    {\choicebitmap[width=0.15\paperwidth]{img/2013/sichuan_science/q7_optC.png}}
    {\choicebitmap[width=0.15\paperwidth]{img/2013/sichuan_science/q7_optD.png}}
\end{problem}

\begin{answer}
C
\end{answer}

\begin{solution}
函数定义域为 \(x\ne0\)，且当 \(x\ne0\) 时 \(\frac{x^3}{3^x-1}>0\). 另外
\[
\lim_{x\to0^-}\frac{x^3}{3^x-1}=\lim_{x\to0^+}\frac{x^3}{3^x-1}=0
\]
但 \(x=0\) 不在定义域内，所以原点是空心点. 当 \(x\to-\infty\) 时函数值趋于 \(+\infty\)，当 \(x\to+\infty\) 时函数值趋于 \(0^+\). 求导得
\[
f'(x)=\frac{x^2\bigl[3(3^x-1)-x3^x\ln3\bigr]}{(3^x-1)^2}
\]
令 \(t=x\ln3\)，记 \(g(t)=3(\e^t-1)-t\e^t\). 当 \(t<0\) 时，\(g'(t)=\e^t(2-t)>0\)，且 \(g(0)=0\)，所以 \(g(t)<0\)，左支递减. 当 \(t>0\) 时，\(g'(t)>0\)（\(0<t<2\)），\(g'(t)<0\)（\(t>2\)），并且 \(g(2)=\e^2-3>0\)、\(g(3)=-3<0\)，故右支先增后减. 因而图象应在原点处有空心点，左右两侧均在 \(x\) 轴上方，只有选项 C 符合，故选 C.
\end{solution}

\begin{problem}
从 \(1\)，\(3\)，\(5\)，\(7\)，\(9\) 这五个数中，每次取出两个不同的数分别为 \(a\)，\(b\)，共可得到 \( \lg a - \lg b \) 的不同值的个数是（\quad）

\choices
    {\(9\)}
    {\(10\)}
    {\(18\)}
    {\(20\)}
\end{problem}

\begin{answer}
C
\end{answer}

\begin{solution}
因为 \(\lg a-\lg b=\lg\frac{a}{b}\)，有序取两个不同的数共有 \(5\times4=20\) 种. 这些比值中只有
\(\frac{1}{3}=\frac{3}{9}\) 和 \(3=\frac{9}{3}\) 各发生一次重复，其他比值不同. 所以不同值的个数为 \(20-2=18\)，故选 C.
\end{solution}

\begin{problem}
节日前夕，小李在家门前的树上挂了两串彩灯，这两串彩灯的第一次闪亮相互独立，且都在通电后的 \(4\text{ 秒}\) 内任一时刻等可能发生，然后每串彩灯以 \(4\text{ 秒}\) 为间隔闪亮. 那么这两串彩灯同时通电后，它们第一次闪亮的时刻相差不超过 \(2\text{ 秒}\) 的概率是（\quad）

\choices
    {\(\frac{1}{4}\)}
    {\( \frac{1}{2} \)}
    {\( \frac{3}{4} \)}
    {\( \frac{7}{8} \)}
\end{problem}

\begin{answer}
C
\end{answer}

\begin{solution}
设两串彩灯第一次闪亮的时刻分别为 \(X\)，\(Y\)，则 \((X,Y)\) 在正方形 \(0\leqslant X\leqslant4\)，\(0\leqslant Y\leqslant4\) 内等可能. 不满足 \(|X-Y|\leqslant2\) 的区域是正方形的两个直角三角形，每个面积为 \(\frac12\cdot2^2=2\)，总面积为 \(4\). 因而
\[
P(|X-Y|\leqslant2)=1-\frac{4}{4^2}=\frac34
\]
故选 C.
\end{solution}

\begin{problem}
设函数 \( f(x) = \sqrt{\e^x + x - a} \) （\( a \in \R \)，\( \e \) 为自然对数的底数）. 若曲线 \( y = \sin x \) 上存在 \( (x_0, y_0) \) 使得 \( f(f(y_0)) = y_0 \)，则 \( a \) 的取值范围是（\quad）

\choices
    {\([1, \e]\)}
    {\([\e^{-1} - 1, 1]\)}
    {\([1, 1 + \e]\)}
    {\([\e^{-1} - 1, \e + 1]\)}
\end{problem}

\begin{answer}
A
\end{answer}

\begin{solution}
令 \(t=y_0=\sin x_0\). 因为 \(f\) 的值非负且 \(f(f(t))=t\)，所以 \(0\leqslant t\leqslant1\). 记 \(u=f(t)\geqslant0\)，则 \(f(u)=t\)，从而
\(u^2=\e^t+t-a\)，\(t^2=\e^u+u-a\).
两式相减并移项得
\[
(u-t)(u+t+1)+(\e^u-\e^t)=0
\]
若 \(u>t\)，左式两项均为正；若 \(u<t\)，左式两项均为负，故只能有 \(u=t\). 于是
\[
a=\e^t+t-t^2
\]
令 \(h(t)=\e^t+t-t^2\)，则在 \(0\leqslant t\leqslant1\) 上
\[
h'(t)=\e^t+1-2t\geqslant(1+t)+1-2t=2-t>0
\]
所以 \(h([0,1])=[h(0),h(1)]=[1,\e]\). 反之，对任意 \(t\in[0,1]\) 都可取 \(x_0\) 使 \(\sin x_0=t\)，且取 \(a=h(t)\) 时 \(f(t)=t\)，故条件确实成立. 因此 \(a\in[1,\e]\)，故选 A.
\end{solution}

\section{填空题}

\begin{problem}
二项式 \((x + y)^5\) 的展开式中，含 \(x^2 y^3\) 的项的系数是\fillinblank{}.
\end{problem}

\begin{answer}
\(10\)
\end{answer}

\begin{solution}
二项式通项为 \(T_{k+1}=\mathrm C_5^k x^{5-k}y^k\). 要得到 \(x^2y^3\)，需 \(5-k=2\)，即 \(k=3\)，所以其系数为
\[
\mathrm C_5^3=\frac{5\cdot4\cdot3}{3\cdot2\cdot1}=10
\]
\end{solution}

\begin{problem}
在平行四边形 \(ABCD\) 中，对角线 \(AC\) 与 \(BD\) 交于点 \(O\)，\(\overrightarrow{AB} + \overrightarrow{AD} = \lambda \overrightarrow{AO}\)，则 \(\lambda =\)\fillinblank{}.
\end{problem}

\begin{answer}
\(2\)
\end{answer}

\begin{solution}
设 \(\overrightarrow{AB}=\bs{b}\)，\(\overrightarrow{AD}=\bs{d}\). 则
\(\overrightarrow{AC}=\bs{b}+\bs{d}\)，而平行四边形的对角线互相平分，所以
\[
\overrightarrow{AO}=\frac12\overrightarrow{AC}=\frac12(\bs{b}+\bs{d})
\]
因此
\[
\overrightarrow{AB}+\overrightarrow{AD}=2\overrightarrow{AO}
\]
故 \(\lambda=2\).
\end{solution}

\begin{problem}
设 \(\sin 2\alpha = -\sin \alpha\)，\(\alpha \in \left(\frac{\pi}{2}, \pi\right)\)，则 \(\tan 2\alpha\) 的值是\fillinblank{}.
\end{problem}

\begin{answer}
\(\sqrt{3}\)
\end{answer}

\begin{solution}
由 \(\sin2\alpha=2\sin\alpha\cos\alpha=-\sin\alpha\)，且
\(\alpha\in(\frac\pi2,\pi)\) 时 \(\sin\alpha\ne0\)，可得
\(2\cos\alpha=-1\)，\(qquad \cos\alpha=-\frac12\).
在给定区间内 \(\alpha=\frac{2\pi}{3}\)，所以
\[
\tan2\alpha=\tan\frac{4\pi}{3}=\sqrt3
\]
\end{solution}

\begin{problem}
已知 \( f(x) \) 是定义域为 \( \R \) 的偶函数，当 \( x \geqslant 0 \) 时，\( f(x) = x^2 - 4x \)，那么，不等式 \( f(x + 2) < 5 \) 的解集是\fillinblank{}.
\end{problem}

\begin{answer}
\((-7,3)\)
\end{answer}

\begin{solution}
由偶性，对任意实数 \(t\) 有 \(f(t)=f(|t|)=|t|^2-4|t|\). 令 \(u=|x+2|\geqslant0\)，则
\[
f(x+2)<5\iff u^2-4u<5\iff (u-5)(u+1)<0
\]
结合 \(u\geqslant0\)，得到 \(0\leqslant u<5\)，即
\[
-5<x+2<5
\]
所以解集为 \((-7,3)\).
\end{solution}

\begin{problem}
设 \(P_1\)，\(P_2\)，\(\cdots\)，\(P_n\) 为平面 \(\alpha\) 内的 \(n\) 个点，在平面 \(\alpha\) 内的所有点中，若点 \(P\) 到 \(P_1\)，\(P_2\)，\(\cdots\)，\(P_n\) 点的距离之和最小，则称点 \(P\) 为 \(P_1\)，\(P_2\)，\(\cdots\)，\(P_n\) 点的一个“中位点”. 例如，线段 \(AB\) 上的任意点都是端点 \(A\)，\(B\) 的中位点. 则有下列命题：

\begin{circlelist}
\item 若三个点 \(A\)，\(B\)，\(C\) 共线，\(C\) 在线段 \(AB\) 上，则 \(C\) 是 \(A\)，\(B\)，\(C\) 的中位点；
\item 直角三角形斜边的中点是该直角三角形三个顶点的中位点；
\item 若四个点 \(A\)，\(B\)，\(C\)，\(D\) 共线，则它们的中位点存在且唯一；
\item 梯形对角线的交点是该梯形四个顶点的唯一中位点.
\end{circlelist}

其中的真命题是\fillinblank{}.（写出所有真命题的序号）
\end{problem}

\begin{answer}
\(\circled{1}\circled{4}\)
\end{answer}

\begin{solution}
对命题 \(\circled{1}\)，任取平面内的点 \(X\)，由三角不等式
\(XA+XB\geqslant AB\)，且 \(XC\geqslant0\). 当 \(X=C\) 时，因 \(C\) 在线段 \(AB\) 上，距离和为 \(CA+CB=AB\)，达到下界，故命题 \(\circled{1}\) 为真.

对命题 \(\circled{2}\)，取边长为 \(3\)，\(4\)，\(5\) 的直角三角形. 斜边中点到三个顶点的距离均为 \(\frac52\)，距离和为 \(\frac{15}{2}\)；直角顶点到三个顶点的距离和为 \(3+4=7<\frac{15}{2}\)，故命题 \(\circled{2}\) 为假.

对命题 \(\circled{3}\)，取同一直线上的四点 \(A\)，\(B\)，\(C\)，\(D\) 的坐标依次为 \(0\)，\(1\)，\(2\)，\(3\). 对任意 \(x\in[1,2]\)，令 \(X\) 为直线上坐标为 \(x\) 的点，则到四点的距离和为
\[
x+(x-1)+(2-x)+(3-x)=4
\]
另一方面，对平面内任意点 \(X\)，由三角不等式有 \(XA+XD\geqslant AD=3\)，\(XB+XC\geqslant BC=1\)，故到四点的距离和不小于 \(4\). 所以线段 \(BC\) 上的点都是中位点，中位点不唯一，命题 \(\circled{3}\) 为假.

对命题 \(\circled{4}\)，设梯形对角线 \(AC\)、\(BD\) 交于 \(O\). 对任意点 \(X\)，有 \(XA+XC\geqslant AC\)，\(XB+XD\geqslant BD\). 当 \(X=O\) 时两式均取等号，所以 \(O\) 的距离和为 \(AC+BD\)，且达到最小值. 若同时取等号，\(X\) 必须同时在线段 \(AC\) 和 \(BD\) 上，故只能是 \(O\)，所以命题 \(\circled{4}\) 为真. 真命题的序号为 \(\circled{1}\)、\(\circled{4}\).
\end{solution}

\section{解答题}

\begin{problem}
在等差数列 \(\{a_{n}\}\) 中，\(a_{1} + a_{3} = 8\)，且 \(a_{4}\) 为 \(a_{2}\) 和 \(a_{9}\) 的等比中项，求数列 \(\{a_{n}\}\) 的首项，公差及前 \(n\) 项和.
\end{problem}

\begin{solution}
设首项为 \(a_1\)，公差为 \(d\)，则 \(a_n=a_1+(n-1)d\). 由
\[
a_1+a_3=2a_1+2d=8
\]
得 \(a_1+d=4\)，即 \(a_2=4\). 又 \(a_4\) 是 \(a_2\)、\(a_9\) 的等比中项，所以
\[
a_4^2=a_2a_9
\]
利用 \(a_1=4-d\)，有 \(a_4=4+2d\)，\(a_9=4+7d\)，从而
\[
(4+2d)^2=4(4+7d)
\iff 4d(d-3)=0
\]
因此 \(d=0\) 或 \(d=3\).

当 \(d=0\) 时，\(a_1=4\)，且 \(S_n=4n\)；当 \(d=3\) 时，\(a_1=1\)，且
\[
S_n=\frac n2\bigl(2a_1+(n-1)d\bigr)=\frac{n(3n-1)}2
\]
所以答案为 \((a_1,d,S_n)=(4,0,4n)\) 或 \((1,3,\frac{n(3n-1)}2)\).
\end{solution}

\begin{problem}
在 \(\triangle ABC\) 中，角 \(A\)，\(B\)，\(C\) 的对边分别为 \(a\)，\(b\)，\(c\)，且 \(2\cos^2\frac{A - B}{2}\cos B - \sin (A - B)\sin B + \cos (A + C) = -\frac{3}{5}\).
\begin{enumerate}
\item 求 \(\cos A\) 的值；
\item 若 \(a = 4\sqrt{2}\)，\(b = 5\)，求向量 \(\overrightarrow{BA}\) 在 \(\overrightarrow{BC}\) 方向上的投影.
\end{enumerate}
\end{problem}

\begin{solution}
\begin{enumerate}
\item 因为 \(A+B+C=\pi\)，且
\[
2\cos^2\frac{A-B}{2}=1+\cos(A-B)
\]
所以题设左端为
\[
\begin{aligned}
&\bigl[1+\cos(A-B)\bigr]\cos B-\sin(A-B)\sin B+\cos(A+C)\\
={}&\cos B+\cos A+\cos(\pi-B)=\cos A.
\end{aligned}
\]
因此 \(\cos A=-\frac35\).

\item 由 \(\cos A=-\frac35\) 及 \(A\in(0,\pi)\)，得 \(\sin A=\frac45\). 由正弦定理
\[
\frac{a}{b}=\frac{\sin A}{\sin B}
\]
及 \(a=4\sqrt2\)、\(b=5\)，得
\[
\sin B=\frac{b\sin A}{a}=\frac{\sqrt2}{2}
\]
因为 \(a>b\)，所以 \(A>B\). 若 \(B=\frac{3\pi}{4}\)，则 \(A>\frac{3\pi}{4}\)，从而 \(\cos A<\cos\frac{3\pi}{4}=-\frac{\sqrt2}{2}\)，这与 \(\cos A=-\frac35>-\frac{\sqrt2}{2}\) 矛盾. 因此 \(B=\frac\pi4\)，即 \(\cos B=\frac{\sqrt2}{2}\).

由余弦定理
\[
a^2=b^2+c^2-2bc\cos A
\]
得
\[
32=25+c^2+6c
\iff c^2+6c-7=0
\]
因 \(c>0\)，故 \(c=1\). 向量 \(\overrightarrow{BA}\) 在 \(\overrightarrow{BC}\) 方向上的投影为
\[
|\overrightarrow{BA}|\cos B=c\cos B=1\cdot\frac{\sqrt2}{2}=\frac{\sqrt2}{2}
\]
\end{enumerate}
\end{solution}

\begin{problem}
某算法的程序框图如图所示，其中输入的变量 \(x\) 在 \(1\)，\(2\)，\(3\)，\(\cdots\)，\(24\) 这 24 个整数中等可能随机产生.

\begin{enumerate}
\item 分别求出按程序框图正确编程运行时输出 \(y\) 的值为 \(i\) 的概率 \(P_{i}\) （\(i = 1\)，\(2\)，\(3\)）；
\item 甲，乙两同学依据自己对程序框图的理解，各自编写程序重复运行 \(n\) 次后，统计记录了输出 \(y\) 的值为 \(i\) （\(i = 1\)，\(2\)，\(3\)） 的频数. 以下是甲，乙所作频数统计表的部分数据. 当 \(n = 2100\) 时，根据表中的数据，分别写出甲，乙所编程序各自输出 \(y\) 的值为 \(i\) （\(i = 1\)，\(2\)，\(3\)） 的频率（用分数表示），并判断两位同学中哪一位所编写程序符合算法要求的可能性较大；

\begin{center}
甲的频数统计表（部分）
\renewcommand{\arraystretch}{1.2}
\begin{tabular}{c|ccc}
\hline
\begin{tabular}{c}运行次数\\\(n\)\end{tabular} & \begin{tabular}{c}输出 \(y\) 的值为\\1 的频数\end{tabular} & \begin{tabular}{c}输出 \(y\) 的值为\\2 的频数\end{tabular} & \begin{tabular}{c}输出 \(y\) 的值为\\3 的频数\end{tabular} \\
\hline
30 & 14 & 6 & 10 \\
\hline
\(\cdots\) & \(\cdots\) & \(\cdots\) & \(\cdots\) \\
\hline
2100 & 1027 & 376 & 697 \\
\hline
\end{tabular}
\end{center}

\begin{center}
乙的频数统计表（部分）
\renewcommand{\arraystretch}{1.2}
\begin{tabular}{c|ccc}
\hline
\begin{tabular}{c}运行次数\\\(n\)\end{tabular} & \begin{tabular}{c}输出 \(y\) 的值为\\1 的频数\end{tabular} & \begin{tabular}{c}输出 \(y\) 的值为\\2 的频数\end{tabular} & \begin{tabular}{c}输出 \(y\) 的值为\\3 的频数\end{tabular} \\
\hline
30 & 12 & 11 & 7 \\
\hline
\(\cdots\) & \(\cdots\) & \(\cdots\) & \(\cdots\) \\
\hline
2100 & 1051 & 696 & 353 \\
\hline
\end{tabular}
\end{center}

\item 按程序框图正确编写的程序运行 3 次，求输出 \(y\) 的值为 2 的次数 \(\xi\) 的分布列及数学期望.
\end{enumerate}

\bitmapfigure[max width=0.48\linewidth]{img/2013/sichuan_science/q18_fig1.png}
\end{problem}

\begin{solution}
\begin{enumerate}
\item \(1\) 至 \(24\) 中有 \(12\) 个奇数，程序输出 \(y=1\)；有 \(12\) 个偶数，其中 \(6\)，\(12\)，\(18\)，\(24\) 共 \(4\) 个能被 \(3\) 整除，输出 \(y=3\)，其余 \(8\) 个输出 \(y=2\). 因而 \(P_1=\frac{12}{24}=\frac12\)，\(P_2=\frac{8}{24}=\frac13\)，\(P_3=\frac{4}{24}=\frac16\).

\item 当 \(n=2100\) 时，甲、乙两程序的频率分别为
\begin{center}
\begin{tabular}{c|ccc}
 & \(y=1\) & \(y=2\) & \(y=3\) \\ \hline
甲 & \(\frac{1027}{2100}\) & \(\frac{376}{2100}\) & \(\frac{697}{2100}\) \\
乙 & \(\frac{1051}{2100}\) & \(\frac{696}{2100}\) & \(\frac{353}{2100}\)
\end{tabular}
\end{center}
正确程序的理论概率为 \(\left(\frac12,\frac13,\frac16\right)\). 以 \(2100\) 为分母，甲的三项频数与理论频数 \(1050\)，\(700\)，\(350\) 分别相差 \(-23\)，\(-324\)，\(347\)，乙分别相差 \(1\)，\(-4\)，\(3\). 因而乙的频率整体更接近理论概率，乙同学所编程序符合算法要求的可能性较大.

\item 每次运行输出 \(y=2\) 的概率为 \(p=\frac13\)，三次运行相互独立. 因而 \(\xi\) 服从二项分布，可能取值为 \(0\)，\(1\)，\(2\)，\(3\)，且
\[
P(\xi=k)=\mathrm C_3^k\left(\frac13\right)^k\left(\frac23\right)^{3-k}
\]
所以分布列为
\begin{center}
\begin{tabular}{c|cccc}
\(\xi\) & 0 & 1 & 2 & 3 \\ \hline
\(P\) & \(\frac{8}{27}\) & \(\frac{4}{9}\) & \(\frac{2}{9}\) & \(\frac{1}{27}\)
\end{tabular}
\end{center}
数学期望为
\[
E(\xi)=3p=3\cdot\frac13=1
\]
\end{enumerate}
\end{solution}

\begin{problem}
如图，在三棱柱 \(ABC - A_{1}B_{1}C_{1}\) 中，侧棱 \(AA_{1} \perp\) 底面 \(ABC\)，\(AB = AC = 2AA_{1}\)，\(\angle BAC = 120^{\circ}\)，\(D\)，\(D_{1}\) 分别是线段 \(BC\)，\(B_{1}C_{1}\) 的中点，\(P\) 是线段 \(AD\) 的中点.
\begin{enumerate}
\item 在平面 \(ABC\) 内，试作出过点 \(P\) 与平面 \(A_{1}BC\) 平行的直线 \(l\)，说明理由，并证明直线 \(l \perp\) 平面 \(ADD_{1}A_{1}\)；
\item 设（1）中的直线 \(l\) 交 \(AB\) 于点 \(M\)，交 \(AC\) 于点 \(N\)，求二面角 \(A-A_{1}M-N\) 的余弦值.
\end{enumerate}
\bitmapfigure[float=inline,max width=0.54\linewidth]{img/2013/sichuan_science/q19_fig1.png}
\end{problem}

\begin{solution}
\begin{enumerate}
\item 在平面 \(ABC\) 内过 \(P\) 作 \(l\parallel BC\). 平面 \(A_1BC\) 与平面 \(ABC\) 的交线为 \(BC\)，而 \(P\notin BC\)，所以 \(l\) 不在平面 \(A_1BC\) 内. 又因 \(BC\subset\) 平面 \(A_1BC\)，由直线与平面平行的判定定理，\(l\parallel\) 平面 \(A_1BC\).

又 \(AB=AC\)，且 \(D\) 是 \(BC\) 的中点，所以 \(AD\perp BC\)，从而 \(AD\perp l\). 在平面 \(ADD_1A_1\) 内过 \(P\) 作直线 \(m\parallel AA_1\)，则 \(m\perp l\)，且 \(m\) 与 \(AD\) 在 \(P\) 相交. 因此 \(l\) 垂直于平面 \(ADD_1A_1\) 内过交点 \(P\) 的两条相交直线 \(AD\)、\(m\)，故 \(l\perp\) 平面 \(ADD_1A_1\).

\item 令 \(AA_1=1\)，建立空间直角坐标系，取
\(A=(0,0,0)\)，\(B=(2,0,0)\)，\(C=(-1,\sqrt3,0)\)，\(A_1=(0,0,1)\).
则 \(D=(\frac12,\frac{\sqrt3}{2},0)\)，\(P=(\frac14,\frac{\sqrt3}{4},0)\). 过 \(P\) 作 \(l\parallel BC\)，由相似关系得 \(M\)，\(N\) 分别为 \(AB\)，\(AC\) 的中点，因此
\(M=(1,0,0)\)，\(N=\left(-\frac12,\frac{\sqrt3}{2},0\right)\).
平面 \(AA_1M\) 的一个法向量为
\[
\bs{n}_1=\overrightarrow{AA_1}\times\overrightarrow{AM}=(0,1,0)
\]
平面 \(A_1MN\) 的一个法向量为
\[
\begin{aligned}
\bs{n}_2
&=\overrightarrow{A_1M}\times\overrightarrow{A_1N}\\
&=(1,0,-1)\times\left(-\frac12,\frac{\sqrt3}{2},-1\right)
=\left(\frac{\sqrt3}{2},\frac32,\frac{\sqrt3}{2}\right).
\end{aligned}
\]
所以二面角的余弦为
\[
\cos\theta=\frac{|\bs{n}_1\cdot\bs{n}_2|}{|\bs{n}_1|\,|\bs{n}_2|}
=\frac{\frac32}{\sqrt{\frac34+\frac94+\frac34}}
=\frac{\sqrt{15}}5
\]
\end{enumerate}
\end{solution}

\begin{problem}
已知椭圆 \(C: \frac{x^2}{a^2} + \frac{y^2}{b^2} = 1\)（\(a>b>0\)）的两个焦点分别为 \(F_1(-1,0)\)，\(F_2(1,0)\)，且椭圆 \(C\) 经过点 \(P\left(\frac{4}{3}, \frac{1}{3}\right)\).
\begin{enumerate}
\item 求椭圆 \(C\) 的离心率；
\item 设过点 \(A(0,2)\) 的直线 \(l\) 与椭圆 \(C\) 交于 \(M\)，\(N\) 两点，点 \(Q\) 是线段 \(MN\) 上的点，且 \(\frac{2}{|AQ|^2} = \frac{1}{|AM|^2} + \frac{1}{|AN|^2}\)，求点 \(Q\) 的轨迹方程.
\end{enumerate}
\end{problem}

\begin{solution}
\begin{enumerate}
\item 椭圆焦半距 \(c=1\). 由椭圆定义，点 \(P\) 到两焦点的距离和为 \(2a\)，而
\(PF_1=\sqrt{\left(\frac43+1\right)^2+\left(\frac13\right)^2}=\frac{5\sqrt2}{3}\)，\(PF_2=\sqrt{\left(\frac43-1\right)^2+\left(\frac13\right)^2}=\frac{\sqrt2}{3}\).
故 \(2a=2\sqrt2\)，\(a=\sqrt2\)，\(b^2=a^2-c^2=1\). 因而离心率
\[
e=\frac ca=\frac{\sqrt2}{2}
\]

\item 由上问，椭圆方程为 \(\frac{x^2}{2}+y^2=1\)，即 \(x^2+2y^2=2\).

先考虑竖直直线 \(x=0\). 此时 \(M=(0,1)\)、\(N=(0,-1)\). 设 \(Q=(0,y)\)，则 \(Q\) 在线段 \(MN\) 上，且
\[
\frac{2}{(2-y)^2}=\frac{1}{1^2}+\frac{1}{3^2}=\frac{10}{9}
\]
因 \(y\leqslant1\)，得 \(y=2-\frac{3\sqrt5}{5}\).

再考虑非竖直直线，设其方程为 \(y=kx+2\). 联立椭圆方程，得交点横坐标 \(x_1\)，\(x_2\) 满足
\[
(1+2k^2)x^2+8kx+6=0
\]
因两交点不同，判别式 \(16k^2-24>0\)，即 \(k^2>\frac32\). 因而
\(x_1+x_2=-\frac{8k}{1+2k^2}\)，\(x_1x_2=\frac{6}{1+2k^2}\)，\(k^2>\frac32\).
设 \(Q=(x,y)\). 因 \(M\)，\(N\)，\(Q\) 在同一直线上，距离均可用横坐标表示，题设条件化为
\[
\frac2{x^2}=\frac1{x_1^2}+\frac1{x_2^2}
\]
于是
\[
x^2=\frac{2x_1^2x_2^2}{x_1^2+x_2^2}
=\frac{18}{10k^2-3}
\]
又 \(k=\frac{y-2}{x}\)，消去 \(k\) 得
\[
10(y-2)^2-3x^2=18
\]
由 \(k^2>\frac32\) 得 \(0<x^2<\frac32\)，即
\[
-\frac{\sqrt6}{2}<x<\frac{\sqrt6}{2}
\]
由于 \(Q\) 在线段 \(MN\) 上，所以 \(y\leqslant1\)，结合轨迹方程只能取下支，从而
\[
\frac12<y<2-\frac{3\sqrt5}{5}
\]
对于非竖直直线，反过来若点 \((x,y)\) 满足上述方程及范围，令 \(k=\frac{y-2}{x}\). 由轨迹方程得 \(x^2=\frac{18}{10k^2-3}\)，且 \(x^2<\frac32\) 从而 \(k^2>\frac32\). 交点横坐标 \(x_1\)，\(x_2\) 同号，且其符号与 \(-k\) 相同；又 \(\operatorname{sgn}x=-\operatorname{sgn}k\)，并且
\[
x^2=\frac{2x_1^2x_2^2}{x_1^2+x_2^2}
\]
是 \(x_1^2\)，\(x_2^2\) 的调和平均数，故 \(x\) 位于 \(x_1\)，\(x_2\) 之间，即 \(Q\) 在线段 \(MN\) 上，且满足题设等式. 最后考虑竖直直线 \(x=0\)：此时 \(M=(0,1)\)、\(N=(0,-1)\)，由题设等式得 \(y=2-\frac{3\sqrt5}{5}\)，正好补上上端点. 因此轨迹为
\(10(y-2)^2-3x^2=18\)，其中 \(-\frac{\sqrt6}{2}<x<\frac{\sqrt6}{2}\)，\(\frac12<y\leqslant2-\frac{3\sqrt5}{5}\).
\end{enumerate}
\end{solution}

\begin{problem}
已知函数 \( f(x) = \begin{cases}
x^2 + 2x + a, & x < 0,\\
\ln x, & x > 0,
\end{cases} \) 其中 \(a\) 是实数. 设 \( A(x_1, f(x_1))\)，\( B(x_2, f(x_2))\) 为该函数图象上的两点，且 \( x_1 < x_2 \).
\begin{enumerate}
\item 指出函数 \( f(x) \) 的单调区间；
\item 若函数 \( f(x) \) 的图象在点 \(A\)，\(B\) 处的切线互相垂直，且 \( x_{2} < 0 \)，求 \( x_{2} - x_{1} \) 的最小值；
\item 若函数 \( f(x) \) 的图象在点 A，B 处的切线重合，求 \(a\) 的取值范围.
\end{enumerate}
\end{problem}

\begin{solution}
\begin{enumerate}
\item 当 \(x<0\) 时，\(f'(x)=2x+2\)，故 \(f\) 在 \((-\infty,-1)\) 上单调递减，在 \([-1,0)\) 上单调递增. 当 \(x>0\) 时，\(f'(x)=\frac1x>0\)，故 \(f\) 在 \((0,+\infty)\) 上单调递增.

\item 因 \(x_2<0\)，两点均在二次函数的左支上. 两处切线斜率分别为 \(2x_1+2\)、\(2x_2+2\)，垂直条件为
\[
(2x_1+2)(2x_2+2)=-1
\]
令 \(u=x_1+1\)，\(v=x_2+1\)，则 \(u<v<1\)，且 \(uv=-\frac14\). 因积为负，\(u<0<v\)，于是 \(v=t\in(0,1)\)，\(u=-\frac1{4t}\). 因而
\[
x_2-x_1=v-u=t+\frac1{4t}\geqslant1
\]
等号在 \(t=\frac12\) 时取得，此时 \(x_2=-\frac12<0\)，符合条件. 所以最小值为 \(1\).

\item 若 \(x_1\)，\(x_2\) 同在负半轴，二次函数两处切线斜率相等将推出 \(x_1=x_2\)；若同在正半轴，对数函数两处切线斜率相等也将推出 \(x_1=x_2\)，均与 \(x_1<x_2\) 矛盾. 因此必有 \(x_1<0<x_2\).

令 \(u=x_1\)、\(v=x_2\). 两条切线重合，先比较斜率得
\(2u+2=\frac1v\)，\(v=\frac1{2(u+1)}\).
因为 \(v>0\)，故 \(-1<u<0\). 再比较截距：二次函数在 \(u\) 处的切线为
\[
y=(2u+2)x+a-u^2
\]
对数函数在 \(v\) 处的切线为
\[
y=\frac{x}{v}+\ln v-1
\]
所以
\[
a=u^2+\ln v-1=u^2-1-\ln\bigl(2(u+1)\bigr)
\]
令 \(t=u+1\in(0,1)\)，则
\[
a=t^2-2t-\ln(2t)=:h(t)
\]
且
\(h'(t)=2t-2-\frac1t=\frac{2t^2-2t-1}{t}<0\)，\((0<t<1)\).
因此 \(h(t)\) 在 \((0,1)\) 上递减，且
\(\lim_{t\to0^+}h(t)=+\infty\)，\(h(1)=-1-\ln2\).
由于 \(t=1\) 不属于定义域对应的 \(u<0\)，下端点不取，故
\[
a\in(-1-\ln2,+\infty)
\]
\end{enumerate}
\end{solution}
